\documentclass{article}

\usepackage{amsmath,amssymb}
\usepackage{xurl}
\usepackage[hidelinks]{hyperref}
\setlength{\emergencystretch}{2em}

% Shared commands for notes in docs/paper-gaps/.

\DeclareUnicodeCharacter{2080}{\ifmmode{}_{0}\else\textsubscript{0}\fi}
\DeclareUnicodeCharacter{2081}{\ifmmode{}_{1}\else\textsubscript{1}\fi}
\DeclareUnicodeCharacter{2082}{\ifmmode{}_{2}\else\textsubscript{2}\fi}
\DeclareUnicodeCharacter{2099}{\ifmmode{}_{n}\else\textsubscript{n}\fi}

\newcommand{\Lean}{\textsc{Lean}}
\ExplSyntaxOn
\NewDocumentCommand{\leanid}{m}{
  \ifmmode
    \textsf{\detokenize{#1}}
  \else
    \group_begin:
    \urlstyle{sf}
    \tl_set:Nn \l_tmpa_tl {#1}
    \tl_replace_all:Nnn \l_tmpa_tl {\_} {_}
    \exp_args:NV \path \l_tmpa_tl
    \group_end:
  \fi
}
\ExplSyntaxOff
\providecommand{\tr}{\operatorname{tr}}
\newcommand{\tnleanIssueBase}{https://github.com/LionSR/TNLean/issues/}
\newcommand{\tnleanPrBase}{https://github.com/LionSR/TNLean/pull/}
\newcommand{\tnleanDocBase}{https://sirui-lu.com/TNLean/docs/}
\newcommand{\ghissue}[1]{\href{\tnleanIssueBase#1}{GitHub issue \##1}}
\newcommand{\ghpr}[1]{\href{\tnleanPrBase#1}{PR \##1}}
\newcommand{\doclink}[1]{\href{\tnleanDocBase#1.html}{\path{#1}}}

% Dirac notation and the identity operator, as in blueprint/src/macros/common.tex.
% \providecommand keeps note-local definitions authoritative.
\usepackage{dsfont}
\providecommand{\ket}[1]{|#1\rangle}
\providecommand{\bra}[1]{\langle #1|}
\providecommand{\braket}[2]{\langle #1 | #2 \rangle}
\providecommand{\ketbra}[2]{|#1\rangle\!\langle #2|}
\providecommand{\Id}{\mathds{1}}
\providecommand{\one}{\mathds{1}}
\providecommand{\C}{\mathbb{C}}
\providecommand{\MN}[1]{M_{#1}(\C)}
\providecommand{\MD}{M_D(\C)}

% Verdict marker: \gapnote{<kind>}{<status>}. Kind is the policy
% classification (clarification, local-correction, scope-restriction,
% unfaithful, false-source, open-gap); status is open, wip (elimination
% underway), resolved, or historical. Machine-read by the site index;
% prints a verdict line.
\newcommand{\gapnote}[2]{\par\noindent\textbf{Verdict:} #1 (#2).\par}


\title{The sweep of the frustration-free area law through range projections}
\author{}
\date{November 20, 2026}
\begin{document}
\maketitle
\gapnote{scope-restriction}{open}

\paragraph{Source.}
The reading note \emph{Entanglement across a cut in two-dimensional
    frustration-free ground states} (November 18, 2026,
\path{docs/audits/2026-11-18_finite_range_frustration_free_area_law.tex},
paragraph ``Contraction in one sweep'', equation \texttt{eq:sweep}) takes
positive contractions $k_i=h_i/J$ with $k_i\Omega=0$ and
$\sum_ik_i\ge g(\Id-P)$, puts $G_i=(\Id-k_i)^{1/2}$ and $D=G_N\cdots G_1$, and
proves
\begin{equation}
    \|D^t(\Id-P)\|\le\theta^t,\qquad \theta=\bigl(1+g/(d+1)^2\bigr)^{-1/2},
    \label{eq:sweep}
\end{equation}
where $d$ bounds the number of supports meeting a given support.

\paragraph{Formal statement.}
The formalization proves \eqref{eq:sweep} only when every $k_i$ is an
orthogonal projection, for which $G_i=\Id-k_i$. It is the detectability lemma
of Anshu, Arad and Vidick in the form already formalized in QICLean, with the
count $d$ of non-commuting partners. The amplification step of the note is
then run with the range projections $\Pi_i$ of the terms $h_i$ in place of the
$k_i$. These satisfy the three properties the note uses of its filters:
\begin{itemize}
    \item $\Pi_i\Omega=0$, since $\Pi_i=h_i^{-1}h_i$ in the functional calculus;
    \item $\Pi_i$ acts on the support of $h_i$;
    \item $\Pi_i\ge h_i/J$ when $0\le h_i\le J\Id$, so that
          $\sum_i\Pi_i\ge g(\Id-P)$ with the same $g=\Delta/J$.
\end{itemize}
Hence the projection case of \eqref{eq:sweep}, applied to the $\Pi_i$, gives
the same contraction factor $\theta$ for a different sweep operator
$D=(\Id-\Pi_N)\cdots(\Id-\Pi_1)$. The light-cone identity
\texttt{eq:cone} and the rest of the amplification argument use only these
three properties, so no hypothesis of the note's model \texttt{eq:model} is
added, and the area law \texttt{eq:goal} is unaffected.

\paragraph{What remains open.}
The note's statement \eqref{eq:sweep} for general positive contractions is not
formalized. It is not needed for \texttt{eq:goal}. Elimination plan: if a later
step needs it, prove the norm-loss argument of the note directly, with
$\|K_iv_N\|\le u_i+\sum_{j>i,\,D_j\cap D_i\ne\varnothing}u_j$ and the
inequality $(1-\sqrt{1-x})^2\le x$ on $[0,1]$.

\end{document}
